\section{Análise dos Paradigmas de Arquitetura (CNN vs. Transformers)}

A comparação entre Redes Neurais Convolucionais (CNN) e modelos baseados em Atenção (Transformers) revelou uma superioridade consistente das arquiteturas convolucionais para o dataset de classificação de algodão. Embora os modelos \textit{Vision Transformers} (ViT) e \textit{Swin Transformer} representem o estado da arte em tarefas de visão computacional generalista em larga escala, seu desempenho neste estudo foi limitado.

Este fenômeno pode ser atribuído à natureza do \textit{inductive bias} (viés indutivo) inerente às CNNs. As convoluções operam sob as premissas de localidade e invariância de translação, o que as torna extremamente eficientes na extração de características de textura e cor em conjuntos de dados de domínio específico e tamanho moderado. Por outro lado, os Transformers possuem uma flexibilidade global que requer volumes massivos de dados para aprender relações espaciais que as CNNs já possuem estruturalmente. No contexto da fibra de algodão, onde os padrões discriminativos são sutis e dependentes de micro-texturas, o reuso de \textit{features} da \textbf{DenseNet} e o refinamento do macro-design da \textbf{ConvNext} mostraram-se paradigmas mais robustos.
